\appendix

\section{Materiales Utilizados}
\label{Anexo2}


\begin{figure}[H]
  \centering
  \begin{subfigure}{0.38\linewidth}
    \fbox{\includegraphics[width=\linewidth]{Alimentacion.png}}
    \caption{Fuente de alimentacion externa}
    \label{fig:alimentacion}
  \end{subfigure}
  \hfill
  \begin{subfigure}{0.42\linewidth}      
    \fbox{\includegraphics[width=\linewidth]{Coco.png}}
    \caption{Cable banana-cocodrilo}
    \label{fig:coco}
  \end{subfigure}
  \begin{subfigure}{0.44\linewidth}      
    \fbox{\includegraphics[width=\linewidth]{kit.png}}
    \caption{Kit de Fibra Óptica. Contiene Transmisor, Receptor, Espejo y Cables}
    \label{fig:kit}
  \end{subfigure}
  \hfill
  \begin{subfigure}{0.44\linewidth}      
    \fbox{\includegraphics[width=\linewidth]{jack.png}}
    \caption{Cable jack 3,5 mm}
    \label{fig:jack}
  \end{subfigure}
  \begin{subfigure}{0.44\linewidth}      
    \fbox{\includegraphics[width=\linewidth]{osc.png}}
    \caption{Osciloscopio Digital Tektronix. Serie TDS.}
    \label{fig:osc}
  \end{subfigure}
  \hfill
  \begin{subfigure}{0.44\linewidth}      
    \fbox{\includegraphics[width=\linewidth]{vfl.png}}
    \caption{Localizador visual de fallas (VFL).}
    \label{fig:vfl}
  \end{subfigure}
  \caption{Materiales Necesarios 1.}
  \label{fig:elementos1}
\end{figure}

\begin{figure}[H]
  \centering
  \begin{subfigure}{0.26\linewidth}      
    \fbox{\includegraphics[width=\linewidth]{medidor.png}}
    \caption{Medidor de Potencia Óptica.}
    \label{fig:medidor}
  \end{subfigure}
  \hfill
  \begin{subfigure}{0.4\linewidth}      
    \fbox{\includegraphics[width=\linewidth]{sc_apc.png}}
    \caption{Conector mecánico SC-APC.}
    \label{fig:sc_apc}
  \end{subfigure}
  \begin{subfigure}{0.44\linewidth}      
    \fbox{\includegraphics[width=\linewidth]{strippers.png}}
    \caption{Pinza Peladora de Precisión.}
    \label{fig:stripper}
  \end{subfigure}
  \hfill
  \begin{subfigure}{0.44\linewidth}      
    \fbox{\includegraphics[width=\linewidth]{cortadora.png}}
    \caption{Cortadora de Precisión para Fibra Óptica.}
    \label{fig:cortadora}
  \end{subfigure}
  \begin{subfigure}{0.48\linewidth}      
    \fbox{\includegraphics[width=\linewidth]{herramienta.png}}
    \caption{Herramienta 3M Fibrlok II Assembly Tool 2501.}
    \label{fig:herramienta}
  \end{subfigure}
  \hfill
  \begin{subfigure}{0.33\linewidth}      
    \fbox{\includegraphics[width=\linewidth]{empalmes.png}}
    \caption{Empalmes mecánicos.}
    \label{fig:empalmes_mecanicos}
  \end{subfigure}
  \begin{subfigure}{0.4\linewidth}      
    \fbox{\includegraphics[width=\linewidth]{protector.png}}
    \caption{Proteccion para empalmes.}
    \label{fig:protector}
  \end{subfigure}
  \caption{Materiales Necesarios 2.}
  \label{fig:elementos2}
\end{figure}

\section{Procedimiento de Conexión de la Alimentación Externa}
\label{Anexo1}

\begin{itemize}
    \item \textbf{Fuente Externa:} Configurada en un rango de $9\,\text{V}$ a $15\,\text{V}$.
    \item \textbf{Conexión Positiva:} Borne positivo de la fuente a los bornes rojos 13 (Transmisor) y 19 (Receptor).
    \item \textbf{Conexión Negativa:} Borne negativo de la fuente a los bornes verdes 15 (Transmisor) y 21 (Receptor).
    \item \textbf{Control de Encendido:} Los interruptores de batería (2 en TX y 9 en RX) no tienen efecto con alimentación externa, por lo que el encendido/apagado se efectúa conectando/desconectando los cables de la fuente.
\end{itemize}